\chapter{Worked Problems}\label{ch:worked}

This chapter puts the book to work. It solves seventeen classic problems from the CSES Problem Set
(\href{https://cses.fi/problemset/}{cses.fi/problemset}) all the way from the statement to accepted code, and it shows the
\emph{thinking} that leads to each solution, not just the answer. The problems follow the order of the book: sorting and
greedy, binary search, prefix sums, graphs, dynamic programming, trees, range queries, strings and binary lifting.

\begin{tip}[Get the most out of this chapter]
Open the problem on CSES, send it to VS Code with Competitive Companion (CPH creates the file from your template) and try it
for twenty minutes \emph{before} reading on. If you are stuck, read only the ``Find the idea'' step and try again. Solving it
yourself is what makes a technique stick; reading a solution only makes it look familiar.
\end{tip}

\section[The method]{The method \hfill \bronze}

Every solution below is found with the same six steps. They are worth making a habit: most wrong answers in contests come
from skipping step 2 (the solution was never fast enough) or step 6 (it was never tested on a hard case).

\begin{figure}[H]
\centering
\begin{tikzpicture}[wbox/.style={draw=cpblue, fill=cpblue!6, rounded corners=3pt, thick, text width=4cm, align=center,
    minimum height=1.3cm, font=\sffamily\small}, node distance=0.6cm]
  \node[wbox] (a) {\textbf{1. Read}\\statement, samples, limits};
  \node[wbox, right=of a] (b) {\textbf{2. Budget}\\limits $\to$ complexity};
  \node[wbox, right=of b] (c) {\textbf{3. Brute force}\\the obvious slow solution};
  \node[wbox, below=of c] (d) {\textbf{4. Observe}\\why is the brute force slow?};
  \node[wbox, left=of d] (e) {\textbf{5. Code}\\template + library};
  \node[wbox, left=of e] (f) {\textbf{6. Test}\\samples, edge cases, stress};
  \draw[ded] (a)--(b); \draw[ded] (b)--(c); \draw[ded] (c)--(d); \draw[ded] (d)--(e); \draw[ded] (e)--(f);
\end{tikzpicture}
\caption{The six steps used for every problem in this chapter.}
\end{figure}

\begin{enumerate}
  \item \textbf{Read} the statement twice. Note what is given, what is asked, and every limit.
  \item \textbf{Budget.} A judge does roughly $10^8$ simple operations per second. Table~\ref{tab:constraints} turns the limits
        into the complexity you need: with $n \le 2\cdot 10^5$, $\Oh{n\log n}$ is fine and $\Oh{n^2} = 4\cdot10^{10}$ is
        hopeless.
  \item \textbf{Brute force.} Write down the simplest correct solution, even if it is far too slow. It fixes what you are
        computing, and later it is the reference for a stress test.
  \item \textbf{Observe.} Ask what the brute force repeats or wastes. The answer is almost always one of the techniques in
        this book, and the ``How to recognize it'' boxes are there to help you spot which.
  \item \textbf{Code} it with your template and, where one exists, a library template pasted with its shortcut.
  \item \textbf{Test}: run the CPH samples, then the extremes (smallest input, largest input, all equal, negative numbers),
        then a stress test against the brute force if anything is uncertain.
\end{enumerate}

\paragraph{About the code.} Each solution is the body of \cc{solve()} in your template: it replaces the empty
\cc|void solve() {}|. Remember what the template gives you: \cc{int} is 64-bit (\cc{#define int long long}), \cc{rep(i, n)}
loops $i = 0, \dots, n-1$ and \cc{rep(i, a, b)} loops $i = a, \dots, b-1$, \cc{per} loops backwards, \cc{cin >> v} reads a
whole vector, \cc{print(a, b)} prints its arguments separated by spaces and ends the line (a vector prints as its elements),
\cc{chmax(a, b)} sets $a = \max(a, b)$, and \cc{INF}, \cc{MOD}, \cc{vi}, \cc{vvi}, \cc{pii}, \cc{pll}, \cc{minpq} are ready
to use. Where a library template is needed (DSU, Fenwick tree) it is pasted above \cc{solve()} with its shortcut.

\begin{fact}[Every program here is tested]
Each solution in this chapter was compiled with your exact template and compiler flags and compared with a brute-force
solution on 150 random inputs (the same kind of stress test as Section~4.4 of GUIDE.pdf). The largest cases were also timed
with the judge's flags: every solution runs in well under 0.25\,s. Statements are paraphrased and limits rounded; always read
the original.
\end{fact}

\begin{center}
\small
\begin{tabularx}{\linewidth}{rlLlr}
\toprule
\thd{\#} & \thd{Problem} & \thd{Technique} & \thd{Chapter} & \thd{Page} \\
\midrule
1 & Maximum Subarray Sum & Kadane's algorithm (a one-variable DP) & \ref{ch:dp} & \pageref{wp:kadane} \\
2 & Apartments & sorting + two pointers, greedy matching & \ref{ch:sorting} & \pageref{wp:apartments} \\
3 & Movie Festival & greedy by earliest end (exchange argument) & \ref{ch:sorting} & \pageref{wp:movies} \\
4 & Factory Machines & binary search on the answer & \ref{ch:binsearch} & \pageref{wp:factory} \\
5 & Subarray Sums II & prefix sums + counting map & \ref{ch:prefix} & \pageref{wp:subarray} \\
6 & Counting Rooms & flood fill, connected components & \ref{ch:graphs} & \pageref{wp:rooms} \\
7 & Message Route & BFS + parent pointers & \ref{ch:graphs} & \pageref{wp:route} \\
8 & Shortest Routes I & Dijkstra & \ref{ch:shortest} & \pageref{wp:dijkstra} \\
9 & Road Construction & DSU (union--find) & \ref{ch:dsu} & \pageref{wp:roads} \\
10 & Coin Combinations II & counting DP (unbounded knapsack) & \ref{ch:dp} & \pageref{wp:coins} \\
11 & Book Shop & 0/1 knapsack & \ref{ch:dp} & \pageref{wp:books} \\
12 & Increasing Subsequence & LIS in $\Oh{n\log n}$ & \ref{ch:dp} & \pageref{wp:lis} \\
13 & Subordinates & tree DP, children before parents & \ref{ch:trees} & \pageref{wp:subordinates} \\
14 & Tree Diameter & two BFS passes & \ref{ch:trees} & \pageref{wp:diameter} \\
15 & Dynamic Range Sum Queries & Fenwick tree & \ref{ch:segtree} & \pageref{wp:bit} \\
16 & String Matching & prefix function (KMP) & \ref{ch:strings} & \pageref{wp:kmp} \\
17 & Planets Queries I & binary lifting on a functional graph & \ref{ch:graphs} & \pageref{wp:planets} \\
\bottomrule
\end{tabularx}
\end{center}

% ======================================================================================================================
\section[Maximum subarray sum]{Maximum subarray sum \hfill \silver}\label{wp:kadane}

\begin{statement}{Maximum Subarray Sum}{CSES Problem Set}
Given an array of $n$ integers, find the largest sum of a contiguous, non-empty subarray.
Limits: $n \le 2\cdot 10^5$, $|a_i| \le 10^9$.
\sample{8\\-1 3 -2 5 3 -5 2 2}{9}
\end{statement}

\wstep{Read the limits.} $n \le 2\cdot10^5$ allows $\Oh{n}$ or $\Oh{n\log n}$. Trying every pair $(l, r)$ is
$\Oh{n^2} = 4\cdot10^{10}$ even with prefix sums. The answer can reach $2\cdot10^5\cdot10^9 = 2\cdot10^{14}$, so it needs
64 bits (your template's \cc{int} already is).

\wstep{Find the idea.} Ask a smaller question: what is the best subarray that \emph{ends} at position $i$? Call its sum
$\mathrm{cur}_i$. Such a subarray is either $a_i$ alone, or $a_i$ appended to the best subarray ending at $i-1$:
\[
  \mathrm{cur}_i = \max(\mathrm{cur}_{i-1} + a_i,\; a_i).
\]
The answer is the largest $\mathrm{cur}_i$. This is a dynamic programming whose whole state is one number, known as
\defn{Kadane's algorithm}. In words: keep a running sum, and start over whenever the sum so far is negative, because a
negative prefix can only make what follows smaller.

\begin{figure}[H]
\centering
\begin{tikzpicture}
  \arr{A}{0}{0}{-1,3,-2,5,3,-5,2,2}\node[lbl, left=4pt of A-1] {$a_i$};
  \arr{C}{0}{-1}{-1,3,1,6,9,4,6,8}\node[lbl, left=4pt of C-1] {cur};
  \arr{B}{0}{-2}{-1,3,3,6,9,9,9,9}\node[lbl, left=4pt of B-1] {best};
  \node[cell, hot] at (C-2.west) {3};
  \node[cell, good] at (C-5.west) {9};
  \node[cell, good] at (B-5.west) {9};
  \node[draw=cpgreen, ultra thick, rounded corners=2pt, fit=(A-2)(A-5), inner sep=1pt] {};
  \draw[brace] (A-2.north west) -- node[above=7pt, lbl] {$3 - 2 + 5 + 3 = 9$} (A-5.north east);
  \node[note, text width=5.6cm, anchor=west] at (6.1, -1)
    {cur is the best sum ending at $i$. Orange: the old sum was negative, so the subarray restarts at $a_i$. Green: the
     best value, reached at the end of the framed subarray.};
\end{tikzpicture}
\caption{Kadane's algorithm on the sample.}
\end{figure}

\wstep{Code.}
\begin{code}
void solve() {
    int n;
    cin >> n;
    ll best = LLONG_MIN, cur = 0;
    rep(i, n) {
        ll x;
        cin >> x;
        cur = max(cur + x, x);  // best sum of a subarray that ends here
        chmax(best, cur);
    }
    print(best);
}
\end{code}

\wstep{Complexity.} $\Oh{n}$ time and $\Oh{1}$ extra memory: the array is never stored.

\begin{pitfall}[Starting from zero]
Initialising \cc{best = 0} is wrong when every number is negative: for $[-3, -1]$ the answer is $-1$, because the subarray
must be non-empty. Start from \cc{LLONG_MIN} (or from $a_0$).
\end{pitfall}

\begin{trick}[Variants you will meet]
\begin{itemize}
  \item \textbf{Length at least $k$}: with prefix sums, the answer is $\max_r \bigl(P_r - \min_{l \le r-k} P_l\bigr)$; keep the
        running minimum.
  \item \textbf{Circular array}: the best wrapping subarray is the total minus the \emph{minimum} subarray. Take the larger
        of that and the ordinary answer, except when all numbers are negative (then the ``wrap'' would be empty).
  \item \textbf{Maximum-sum rectangle} in an $n\times m$ matrix: fix the top and bottom rows, add the columns between them
        into one array, and run Kadane on it: $\Oh{n^2 m}$.
\end{itemize}
\end{trick}

% ======================================================================================================================
\section[Apartments]{Apartments \hfill \silver}\label{wp:apartments}

\begin{statement}{Apartments}{CSES Problem Set}
There are $n$ applicants and $m$ apartments. Applicant $i$ wants an apartment of size $a_i$ and accepts any size in
$[a_i - k, a_i + k]$. Each apartment can go to at most one applicant. How many applicants can get an apartment?
Limits: $n, m \le 2\cdot 10^5$; sizes and $k$ up to $10^9$.
\sample{4 3 5\\60 45 80 60\\30 60 75}{2}
\end{statement}

\wstep{Read the limits.} Pairing every applicant with every apartment is $4\cdot10^{10}$ pairs; general matching is much
worse. We need sorting and a single sweep: $\Oh{n\log n + m\log m}$.

\wstep{Find the idea.} Sort both lists and look at the smallest apartment $b_j$ and the smallest applicant $a_i$:
\begin{itemize}
  \item If $b_j < a_i - k$, the apartment is too small for applicant $i$, and for everyone after $i$ too (they want even
        more). Nobody can use it: skip the apartment.
  \item If $b_j > a_i + k$, the apartment is too big for $i$, and every remaining apartment is even bigger. Applicant $i$
        can never be served: skip the applicant.
  \item Otherwise $b_j$ fits $a_i$: match them. Giving the smallest usable apartment to the smallest applicant never hurts,
        because any later applicant who could use $b_j$ could also use the larger apartments that $a_i$ would otherwise
        take (an exchange argument, Section~\ref{sec:exchange}).
\end{itemize}
Every step discards an applicant or an apartment, so there are at most $n + m$ steps: this is the \defn{two pointers}
technique.

\begin{figure}[H]
\centering
\begin{tikzpicture}
  \draw[->, cpgray] (0,0) -- (8.6,0) node[right, lbl] {size};
  \foreach \av/\y/\st in {45/0.7/dim, 60/1.2/good, 60/1.7/dim, 80/2.2/good} {
    \node[\st, draw, rounded corners=2pt, minimum height=3.5mm, minimum width=1.3cm, inner sep=0pt, anchor=west]
      at ({(\av-30)*0.13}, \y) {};
    \node[lbl, anchor=east] at ({(\av-30)*0.13 - 0.08}, \y) {$a = \av$};
  }
  \draw[cpgreen, thick, dashed] (4.55,0.1) -- (4.55,1.02);
  \draw[cpgreen, thick, dashed] (6.5,0.1) -- (6.5,2.02);
  \foreach \bv/\st in {30/dim, 60/good, 75/good} {
    \node[circle, \st, draw, thick, minimum size=2.4mm, inner sep=0pt] (b\bv) at ({(\bv-25)*0.13}, 0) {};
    \node[lbl, below=2pt of b\bv] {\bv};
  }
\end{tikzpicture}
\caption{Sorted applicants (bars: acceptable sizes $[a_i - k, a_i + k]$ with $k = 5$) and apartments (dots). The sweep
matches 60 with the first applicant who wants 60 and 75 with the applicant who wants 80; apartment 30 fits nobody.}
\end{figure}

\wstep{Code.}
\begin{code}
void solve() {
    int n, m, k;
    cin >> n >> m >> k;
    vi a(n), b(m);
    cin >> a >> b;
    sort(all(a)), sort(all(b));
    int ans = 0;
    for (int i = 0, j = 0; i < n && j < m;) {
        if (b[j] < a[i] - k) j++;       // too small for everyone still waiting
        else if (b[j] > a[i] + k) i++;  // no apartment left fits applicant i
        else ans++, i++, j++;
    }
    print(ans);
}
\end{code}

\wstep{Complexity.} $\Oh{n\log n + m\log m}$ for sorting; the sweep is $\Oh{n + m}$.

\begin{recognize}
Two sets must be paired, and whether $x$ and $y$ can be paired depends on comparing their values (a threshold or a
window). Sort both sides and sweep them together; the greedy match of ``smallest with smallest'' is usually optimal.
\end{recognize}
\usetemplate{Two pointers (sliding window)}{Ctrl+Alt+T, 5}{cptwoptr}

% ======================================================================================================================
\section[Movie festival]{Movie festival \hfill \silver}\label{wp:movies}

\begin{statement}{Movie Festival}{CSES Problem Set}
There are $n$ movies; movie $i$ runs from time $a_i$ to $b_i$ ($a_i < b_i$). You can watch one movie at a time and only
from start to end; a movie may start exactly when the previous one ends. What is the largest number of whole movies you
can watch? Limits: $n \le 2\cdot 10^5$, times up to $10^9$.
\sample{3\\3 5\\4 9\\5 8}{2}
\end{statement}

\wstep{Read the limits.} Trying all $2^n$ subsets works for $n \le 20$; here $n$ is $2\cdot10^5$, so we need a greedy rule
and $\Oh{n\log n}$.

\wstep{Find the idea.} Which movie should you watch first? The one that \emph{ends} earliest, because it leaves the most
time for everything else. Proof by exchange: take any optimal schedule and replace its first movie by the movie that ends
earliest overall. That movie ends no later, so it does not clash with the rest of the schedule, and the count is the same.
So some optimal schedule starts with the earliest-ending movie; apply the same argument to the movies that start after it
ends. In code: sort by end time and take every movie that starts no earlier than the end of the last one taken.

\begin{figure}[H]
\centering
\begin{tikzpicture}[x=0.9cm]
  \draw[->, cpgray] (0,0) -- (10.5,0) node[right, lbl] {time};
  \foreach \t in {0,...,10} \draw[cpgray] (\t,0.07) -- (\t,-0.07) node[below, lbl] {\t};
  \node[seg, good, minimum width=1.8cm, anchor=west] at (3,0.6) {3--5};
  \node[seg, good, minimum width=2.7cm, anchor=west] at (5,1.25) {5--8};
  \node[seg, dim, minimum width=4.5cm, anchor=west] at (4,1.9) {4--9};
  \node[lbl, anchor=west] at (10.8,0.6) {ends first: take it};
  \node[lbl, anchor=west] at (10.8,1.25) {starts at $5 \ge 5$: take it};
  \node[lbl, anchor=west] at (10.8,1.9) {starts at $4 < 8$: skip};
\end{tikzpicture}
\caption{Movies sorted by end time; the greedy takes the green ones.}
\end{figure}

\wstep{Code.}
\begin{code}
void solve() {
    int n;
    cin >> n;
    vector<pii> mv(n);
    for (auto& [a, b] : mv) cin >> a >> b;
    sort(all(mv), [](const pii& x, const pii& y) { return x.second < y.second; });
    int cnt = 0, free_at = 0;  // the time the last chosen movie ends
    for (auto [a, b] : mv)
        if (a >= free_at) cnt++, free_at = b;
    print(cnt);
}
\end{code}

\wstep{Complexity.} $\Oh{n\log n}$ for the sort, then one pass.

\begin{pitfall}[Plausible rules that are wrong]
\begin{itemize}
  \item \emph{Earliest start first}: for $[1,10], [2,3], [4,5]$ it takes $[1,10]$ and nothing else, but the answer is 2.
  \item \emph{Shortest first}: for $[1,5], [4,7], [6,10]$ it takes $[4,7]$, which blocks both others; the answer is 2.
\end{itemize}
Whenever you invent a greedy rule, look for a small counterexample (or stress test it) before coding it.
\end{pitfall}
\usetemplate{Intervals: merge + max disjoint}{Ctrl+Alt+T, Shift+N}{cpintervals}

% ======================================================================================================================
\section[Factory machines]{Factory machines \hfill \silver}\label{wp:factory}

\begin{statement}{Factory Machines}{CSES Problem Set}
A factory has $n$ machines working at the same time; machine $i$ makes one product every $k_i$ seconds. What is the
shortest time in which the machines can make $t$ products together? Limits: $n \le 2\cdot 10^5$; $t, k_i \le 10^9$.
\sample{3 7\\3 2 5}{8}
\end{statement}

\wstep{Read the limits.} The answer can be as large as $10^{18}$ (one machine with $k = t = 10^9$), so trying times one
second at a time is out of the question.

\wstep{Find the idea.} Turn ``find the minimum time'' into a yes/no question: \emph{can the machines make $t$ products within
time $T$?} That is easy: machine $i$ makes $\lfloor T / k_i\rfloor$ products, so add these up and compare with $t$, in
$\Oh{n}$. The answer is also \emph{monotone}: if time $T$ is enough, so is $T+1$. A monotone yes/no question means
\defn{binary search on the answer} (Chapter~\ref{ch:binsearch}): find the first $T$ that says yes. The search range is
$[1, \min_i k_i \cdot t]$, because the fastest machine alone makes $t$ products in that time.

\begin{figure}[H]
\centering
\begin{tikzpicture}
  \arr{M}{0}{0}{0,1,2,3,4,6,6,7,8,10}
  \foreach \i in {1,...,10} \node[lbl, above=0pt of M-\i] {\i};
  \node[lbl, left=4pt of M-1] {made};
  \node[lbl, anchor=east] at ($(M-1.north west)+(-0.1,0.2)$) {$T$};
  \arr{V}{0}{-0.9}{no,no,no,no,no,no,no,yes,yes,yes}
  \foreach \i in {1,...,7} \node[cell, fill=cpred!12, draw=cpred!60] at (V-\i.west) {no};
  \foreach \i in {8,...,10} \node[cell, good] at (V-\i.west) {yes};
  \node[lbl, left=4pt of V-1] {$\ge 7$?};
  \draw[hded] ($(V-8.south)+(0,-0.75)$) node[below, lbl] {first yes: $T = 8$} -- (V-8.south);
  \node[note, text width=4.6cm, anchor=west] at (7.4, -0.45)
    {made$(T) = \sum_i \lfloor T / k_i \rfloor$ for $k = (3, 2, 5)$. The answers form a block of ``no'' followed by a
     block of ``yes'', so binary search finds the boundary.};
\end{tikzpicture}
\caption{Products made by time $T$ in the sample, and the monotone yes/no question.}
\end{figure}

\wstep{Code.}
\begin{code}
void solve() {
    int n, t;
    cin >> n >> t;
    vi k(n);
    cin >> k;
    auto ok = [&](ll T) {  // can the machines make t products within time T?
        ll made = 0;
        for (ll x : k)
            if ((made += T / x) >= t) return true;
        return false;
    };
    ll lo = 1, hi = *min_element(all(k)) * t;  // the fastest machine alone is enough
    while (lo < hi) {
        ll mid = lo + (hi - lo) / 2;
        if (ok(mid)) hi = mid;
        else lo = mid + 1;
    }
    print(lo);
}
\end{code}

\wstep{Complexity.} $\Oh{n \log(\min k \cdot t)}$: about 60 checks of $\Oh{n}$ each, roughly $1.2\cdot10^7$ operations.

\begin{pitfall}[Overflow inside the check]
Summing $\lfloor T/k_i\rfloor$ over all machines can reach $2\cdot10^5 \cdot 10^{18}$, far beyond 64 bits. The check stops
as soon as the count reaches $t$: before each addition the count is below $t \le 10^9$, so it never exceeds about
$10^{18} + 10^9$, which fits. Also write the midpoint as \cc{lo + (hi - lo) / 2}.
\end{pitfall}
\usetemplate{Binary search on the answer}{Ctrl+Alt+T, A}{cpbsearch}

% ======================================================================================================================
\section[Subarray sums]{Subarray sums \hfill \silver}\label{wp:subarray}

\begin{statement}{Subarray Sums II}{CSES Problem Set}
Given $n$ integers (which may be negative) and a target $x$, count the contiguous subarrays whose sum is exactly $x$.
Limits: $n \le 2\cdot 10^5$, $|a_i|, |x| \le 10^9$.
\sample{5 7\\2 -1 3 5 -2}{2}
\end{statement}

\wstep{Read the limits.} All $\Oh{n^2}$ pairs are too many. With only positive numbers a sliding window would work (the sum
grows as the window grows), but negative numbers break that: this is exactly what separates ``Subarray Sums II'' from
``Subarray Sums I''.

\wstep{Find the idea.} With prefix sums $P_0 = 0$ and $P_r = a_1 + \dots + a_r$, the sum of $a_{l+1}, \dots, a_r$ is
$P_r - P_l$ (Chapter~\ref{ch:prefix}). So we count pairs $l < r$ with $P_l = P_r - x$. Sweep $r$ from left to right and keep,
in a map, how many earlier prefixes had each value: at each $r$, add the number of earlier prefixes equal to $P_r - x$, then
record $P_r$.

\begin{figure}[H]
\centering
\begin{tikzpicture}
  \arr{P}{0}{0}{0,2,1,4,9,7}\idx{P}{6}
  \node[lbl, left=4pt of P-1] {$P$};
  \arr{A}{0.35}{-1.35}{2,-1,3,5,-2}
  \node[lbl, left=4pt of A-1] {$a$};
  \draw[hded] (P-2.north) to[out=60, in=120] node[above, lbl, text=cporange] {$9 - 2 = 7$} (P-5.north);
  \draw[ultra thick, cpteal, ->] (P-1.north) to[out=75, in=105, looseness=1.3] node[above, lbl, text=cpteal] {$7 - 0 = 7$}
    (P-6.north);
  \node[note, text width=6cm, anchor=west] at (4.8, -0.4)
    {$P_r - P_l = x$ exactly when $a_{l+1} + \dots + a_r = x$. Orange: $P_4 - P_1$, the subarray $-1, 3, 5$. Teal:
     $P_5 - P_0$, the whole array (this one needs the empty prefix $P_0 = 0$).};
\end{tikzpicture}
\caption{The two subarrays with sum 7 in the sample, seen as pairs of prefix sums that differ by 7.}
\end{figure}

\wstep{Code.}
\begin{code}
void solve() {
    int n;
    ll x;
    cin >> n >> x;
    map<ll, int> seen;  // how many prefixes had each sum
    seen[0] = 1;
    ll pre = 0, ans = 0;
    rep(i, n) {
        ll a;
        cin >> a;
        pre += a;
        auto it = seen.find(pre - x);  // earlier prefix P with pre - P = x
        if (it != seen.end()) ans += it->second;
        seen[pre]++;
    }
    print(ans);
}
\end{code}

\wstep{Complexity.} $\Oh{n\log n}$ with \cc{map}. The answer can reach $n(n+1)/2 \approx 2\cdot10^{10}$, which needs 64 bits.

\begin{pitfall}[The empty prefix]
\cc{seen[0] = 1} stands for $P_0$, the sum of nothing. Without it every subarray that starts at the first element is
missed; in the sample the answer would be 1.
\end{pitfall}

\begin{tip}
\cc{unordered_map} is faster on average, but on Codeforces it can be attacked with inputs that make every key collide; use
\cc{map}, or a hash map with a randomized hash (Chapter~\ref{ch:ds-basics}).
\end{tip}
\usetemplate{Prefix sums (1D)}{Ctrl+Alt+T, 1}{cpprefix}

% ======================================================================================================================
\section[Counting rooms]{Counting rooms \hfill \silver}\label{wp:rooms}

\begin{statement}{Counting Rooms}{CSES Problem Set}
A building map has $n$ rows and $m$ columns; each cell is floor (\texttt{.}) or wall (\texttt{\#}). You can walk left,
right, up and down through floor cells. Count the rooms. Limits: $n, m \le 1000$.
\sample{5 8\\\#\#\#\#\#\#\#\#\\\#..\#...\#\\\#\#\#\#.\#.\#\\\#..\#...\#\\\#\#\#\#\#\#\#\#}{3}
\end{statement}

\wstep{Read the limits.} Up to $10^6$ cells, so every cell may be touched only a constant number of times: $\Oh{nm}$.

\wstep{Find the idea.} A grid is a graph in disguise: every floor cell is a vertex, and neighbouring floor cells are joined by
an edge. A room is a \defn{connected component}. Scan the cells row by row; the first time you meet an unvisited floor cell
you have found a new room, so count it and \defn{flood fill} it (a BFS that marks every cell it reaches), so that none of
its cells starts another count. Overwriting visited cells with \texttt{\#} saves a separate visited array.

\begin{figure}[H]
\centering
\begin{tikzpicture}
  \foreach \r in {0,...,4} \foreach \c in {0,...,7} \fill[cpgray!45] (\c*0.55, -\r*0.55) rectangle ++(0.55, -0.55);
  \foreach \r/\c in {1/1, 1/2} \fill[cporange!35] (\c*0.55, -\r*0.55) rectangle ++(0.55, -0.55);
  \foreach \r/\c in {1/4, 1/5, 1/6, 2/4, 2/6, 3/4, 3/5, 3/6} \fill[cpteal!30] (\c*0.55, -\r*0.55) rectangle ++(0.55, -0.55);
  \foreach \r/\c in {3/1, 3/2} \fill[cppurple!25] (\c*0.55, -\r*0.55) rectangle ++(0.55, -0.55);
  \draw[white, thick] (0,-2.75) grid[step=0.55] (4.4,0);
  \foreach \r/\c/\rid in {1/1/1, 1/4/2, 3/1/3} \node[font=\sffamily\scriptsize\bfseries] at ({(\c+0.5)*0.55}, {-(\r+0.5)*0.55}) {\rid};
  \node[note, text width=6.4cm, anchor=west] at (5, -1.35)
    {Walls are grey. Scanning row by row, the numbered cell is the first free cell met in each room. Each one starts a flood
     fill that paints its whole room, so later cells of that room are skipped and each room is counted exactly once.};
\end{tikzpicture}
\caption{The three rooms of the sample.}
\end{figure}

\wstep{Code.}
\begin{code}
void solve() {
    int n, m;
    cin >> n >> m;
    vector<string> g(n);
    cin >> g;
    const int dr[4] = {-1, 1, 0, 0}, dc[4] = {0, 0, -1, 1};
    int rooms = 0;
    rep(sr, n) rep(sc, m) {
        if (g[sr][sc] != '.') continue;
        rooms++;  // a new room: flood it so it is not counted again
        queue<pii> q;
        q.push({sr, sc});
        g[sr][sc] = '#';
        while (!q.empty()) {
            auto [r, c] = q.front();
            q.pop();
            rep(d, 4) {
                int nr = r + dr[d], nc = c + dc[d];
                bool inside = 0 <= nr && nr < n && 0 <= nc && nc < m;
                if (!inside || g[nr][nc] != '.') continue;
                g[nr][nc] = '#';
                q.push({nr, nc});
            }
        }
    }
    print(rooms);
}
\end{code}

\wstep{Complexity.} $\Oh{nm}$: every cell enters the queue at most once and looks at four neighbours.

\begin{pitfall}[Recursive flood fill]
A recursive DFS can go $10^6$ calls deep on a $1000\times1000$ grid. Your local build has a 256\,MB stack, so it passes on
your machine, but on judges with the usual 1--8\,MB stack it crashes. The queue-based BFS above has no such limit.
\end{pitfall}
\usetemplate{Flood fill (grid regions)}{Ctrl+Alt+T, F}{cpflood}

% ======================================================================================================================
\section[Message route]{Message route \hfill \silver}\label{wp:route}

\begin{statement}{Message Route}{CSES Problem Set}
A network has $n$ computers and $m$ two-way connections. Find a route from computer 1 to computer $n$ that passes through as
few computers as possible: print their number and the route, or \texttt{IMPOSSIBLE}. Limits: $n \le 10^5$, $m \le 2\cdot10^5$.
\sample{5 5\\1 2\\1 3\\1 4\\2 3\\5 4}{3\\1 4 5}
\end{statement}

\wstep{Read the limits.} $\Oh{n + m}$ or close to it.

\wstep{Find the idea.} All connections cost the same, so the fewest computers means the fewest edges: an unweighted shortest
path, which \defn{BFS} finds (Chapter~\ref{ch:graphs}). BFS discovers vertices in order of distance. To print the route,
remember for every vertex the vertex it was discovered from, \cc{par[v]}. Following \cc{par} back from $n$ walks the
shortest path backwards; reverse it.

\begin{figure}[H]
\centering
\begin{tikzpicture}
  \node[vtx, good] (1) at (0,0) {1};
  \node[vtx] (2) at (1.9,1.25) {2};
  \node[vtx] (3) at (1.9,0) {3};
  \node[vtx, good] (4) at (1.9,-1.25) {4};
  \node[vtx, good] (5) at (3.8,-1.25) {5};
  \draw[ed, dashed] (2) -- (3);
  \draw[hded] (1) -- (2); \draw[hded] (1) -- (3);
  \draw[ultra thick, cpgreen, ->] (1) -- (4); \draw[ultra thick, cpgreen, ->] (4) -- (5);
  \node[lbl, left=2pt of 1] {$d = 0$}; \node[lbl, above=2pt of 2] {$d = 1$}; \node[lbl, right=2pt of 3] {$d = 1$};
  \node[lbl, below=2pt of 4] {$d = 1$}; \node[lbl, below=2pt of 5] {$d = 2$};
  \node[note, text width=6.2cm, anchor=west] at (5, 0)
    {Arrows point from \texttt{par[v]} to $v$: the edge through which $v$ was discovered. Walking back from 5 gives $5, 4, 1$;
     reversed, the route is $1\ 4\ 5$. The dashed edge joins two vertices of the same layer and is never used.};
\end{tikzpicture}
\caption{BFS from computer 1 on the sample, with parent pointers.}
\end{figure}

\wstep{Code.}
\begin{code}
void solve() {
    int n, m;
    cin >> n >> m;
    vvi adj(n);
    rep(i, m) {
        int u, v;
        cin >> u >> v;
        u--, v--;
        adj[u].push_back(v), adj[v].push_back(u);
    }
    vi dist(n, -1), par(n, -1);
    queue<int> q;
    q.push(0), dist[0] = 0;
    while (!q.empty()) {
        int u = q.front();
        q.pop();
        for (int v : adj[u])
            if (dist[v] == -1) dist[v] = dist[u] + 1, par[v] = u, q.push(v);
    }
    if (dist[n - 1] == -1) return print("IMPOSSIBLE");
    vi path;
    for (int v = n - 1; v != -1; v = par[v]) path.push_back(v + 1);
    reverse(all(path));
    print(sz(path));
    print(path);
}
\end{code}

\wstep{Complexity.} $\Oh{n + m}$.

\begin{pitfall}
\begin{itemize}
  \item DFS does \emph{not} find shortest paths; only BFS (or Dijkstra with weights) does.
  \item Mark a vertex when it is \emph{pushed}, not when it is popped. Otherwise a vertex can enter the queue once per
        neighbour, and the queue can grow to $\Oh{m}$ entries with repeated work.
\end{itemize}
\end{pitfall}

\begin{tip}
\cc{return print("IMPOSSIBLE");} prints and leaves \cc{solve()} in one line. It is legal because both functions return
\cc{void}; it saves an \cc{else} and a level of indentation.
\end{tip}
\usetemplate{BFS (shortest path, unweighted)}{Ctrl+Alt+T, B}{cpbfs}

% ======================================================================================================================
\section[Shortest routes]{Shortest routes \hfill \gold}\label{wp:dijkstra}

\begin{statement}{Shortest Routes I}{CSES Problem Set}
There are $n$ cities and $m$ one-way flights; each flight has a price. Print the cheapest total price from city 1 to every
city (all cities can be reached). Limits: $n \le 10^5$, $m \le 2\cdot10^5$, prices up to $10^9$.
\sample{3 4\\1 2 6\\1 3 2\\3 2 3\\1 3 4}{0 5 2}
\end{statement}

\wstep{Read the limits.} Bellman--Ford's $\Oh{nm} = 2\cdot10^{10}$ is too slow. Weights are non-negative, so
\defn{Dijkstra} in $\Oh{(n+m)\log m}$ applies (Chapter~\ref{ch:shortest}).

\wstep{Find the idea.} Always settle the unsettled city with the smallest known distance: with non-negative weights no later
path can make it cheaper. A min-priority queue of (distance, city) pairs gives that city in $\Oh{\log m}$. When a flight
improves a distance, push the new pair and leave the old one in the queue; when an outdated pair is popped
(\cc{d != dist[u]}), skip it. This \emph{lazy deletion} is simpler than updating entries inside the heap.

\begin{figure}[H]
\centering
\begin{minipage}{0.36\linewidth}
\centering
\begin{tikzpicture}
  \node[vtx, good] (1) at (0,0) {1};
  \node[vtx] (3) at (1.8,-1.4) {3};
  \node[vtx] (2) at (3.6,0) {2};
  \draw[ded] (1) -- node[above, lbl] {6} (2);
  \draw[hded] (1) to[bend right=22] node[below left, lbl] {2} (3);
  \draw[ded] (1) to[bend left=22] node[above right, lbl] {4} (3);
  \draw[hded] (3) -- node[below right, lbl] {3} (2);
  \node[lbl, above=2pt of 1] {dist 0}; \node[lbl, above=2pt of 2] {dist 5}; \node[lbl, below=2pt of 3] {dist 2};
\end{tikzpicture}
\end{minipage}\hfill
\begin{minipage}{0.6\linewidth}
\small
\begin{tabular}{@{}clp{5.4cm}@{}}
\toprule
\thd{Step} & \thd{Pop} & \thd{What happens} \\
\midrule
1 & $(0, 1)$ & dist$[2] = 6$, dist$[3] = 2$; the second flight $1\to3$ (price 4) is no better \\
2 & $(2, 3)$ & $2 + 3 = 5 < 6$: dist$[2] = 5$, push $(5, 2)$ \\
3 & $(5, 2)$ & city 2 has no flights out \\
4 & $(6, 2)$ & outdated: $6 \ne \mathrm{dist}[2] = 5$, skipped \\
\bottomrule
\end{tabular}
\end{minipage}
\caption{Dijkstra on the sample. Orange: the cheapest routes.}
\end{figure}

\wstep{Code.}
\begin{code}
void solve() {
    int n, m;
    cin >> n >> m;
    vector<vector<pll>> adj(n);
    rep(i, m) {
        int a, b;
        ll c;
        cin >> a >> b >> c;
        adj[a - 1].push_back({b - 1, c});  // one-way flight
    }
    vector<ll> dist(n, INF);
    minpq<pll> pq;
    dist[0] = 0, pq.push({0, 0});
    while (!pq.empty()) {
        auto [d, u] = pq.top();
        pq.pop();
        if (d != dist[u]) continue;
        for (auto [v, w] : adj[u])
            if (d + w < dist[v]) dist[v] = d + w, pq.push({dist[v], v});
    }
    print(dist);
}
\end{code}

\wstep{Complexity.} $\Oh{(n + m)\log m}$; 0.14\,s on a random graph of the maximum size.

\begin{pitfall}
\begin{itemize}
  \item Distances reach $(n-1)\cdot10^9 \approx 10^{14}$: 64-bit integers are required.
  \item Without the \cc{d != dist[u]} check the answer is still right, but a city can be expanded once per outdated entry,
        which is quadratic on some graphs (a star with many improving edges).
  \item With \emph{negative} weights Dijkstra is wrong; use Bellman--Ford or SPFA with care.
\end{itemize}
\end{pitfall}
\usetemplate{Dijkstra}{Ctrl+Alt+T, J}{cpdijkstra}

% ======================================================================================================================
\section[Road construction]{Road construction \hfill \gold}\label{wp:roads}

\begin{statement}{Road Construction}{CSES Problem Set}
There are $n$ cities and no roads. Then $m$ two-way roads are built one by one. After each road, print the number of
connected components and the size of the largest one. Limits: $n \le 10^5$, $m \le 2\cdot10^5$.
\sample{5 3\\1 2\\1 3\\4 5}{4 2\\3 3\\2 3}
\end{statement}

\wstep{Read the limits.} A BFS after every road costs $\Oh{m(n+m)} \approx 6\cdot10^{10}$. Each road must be handled in
(almost) constant time.

\wstep{Find the idea.} Components only ever \emph{merge}, never split: exactly what a \defn{disjoint set union} (DSU) is made
for (Chapter~\ref{ch:dsu}). \cc{unite(a, b)} returns \cc{true} if it merged two different components; the DSU also keeps
the number of components and each component's size. The largest component can only grow, so a running maximum updated after
each successful merge is enough.

\begin{figure}[H]
\centering
\begin{tikzpicture}
  \begin{scope}
    \node[vtx, hot] (1) at (0,1) {1}; \node[vtx, hot] (2) at (1.1,1) {2}; \node[vtx] (3) at (0.55,0) {3};
    \node[vtx] (4) at (2.4,1) {4}; \node[vtx] (5) at (2.4,0) {5};
    \draw[hed] (1) -- (2);
    \node[lbl] at (1.2,-0.75) {after road 1--2}; \node[font=\ttfamily\small] at (1.2,-1.2) {4 2};
  \end{scope}
  \begin{scope}[xshift=4.4cm]
    \node[vtx, hot] (1) at (0,1) {1}; \node[vtx, hot] (2) at (1.1,1) {2}; \node[vtx, hot] (3) at (0.55,0) {3};
    \node[vtx] (4) at (2.4,1) {4}; \node[vtx] (5) at (2.4,0) {5};
    \draw[ed] (1) -- (2); \draw[hed] (1) -- (3);
    \node[lbl] at (1.2,-0.75) {after road 1--3}; \node[font=\ttfamily\small] at (1.2,-1.2) {3 3};
  \end{scope}
  \begin{scope}[xshift=8.8cm]
    \node[vtx, hot] (1) at (0,1) {1}; \node[vtx, hot] (2) at (1.1,1) {2}; \node[vtx, hot] (3) at (0.55,0) {3};
    \node[vtx, cool] (4) at (2.4,1) {4}; \node[vtx, cool] (5) at (2.4,0) {5};
    \draw[ed] (1) -- (2); \draw[ed] (1) -- (3); \draw[hed] (4) -- (5);
    \node[lbl] at (1.2,-0.75) {after road 4--5}; \node[font=\ttfamily\small] at (1.2,-1.2) {2 3};
  \end{scope}
\end{tikzpicture}
\caption{Components after each road (one colour per component of size at least 2) and the printed line.}
\end{figure}

\wstep{Code.} Paste the DSU template above \cc{solve()} with \key{Ctrl+Alt+T, U}, then:
\begin{code}
void solve() {  // DSU (cpdsu) is pasted above solve()
    int n, m;
    cin >> n >> m;
    DSU d(n);
    int biggest = 1;
    rep(i, m) {
        int a, b;
        cin >> a >> b;
        if (d.unite(a - 1, b - 1)) chmax(biggest, d.size(a - 1));
        print(d.comps, biggest);
    }
}
\end{code}

\wstep{Complexity.} $\Oh{m\,\alpha(n)}$, effectively $\Oh{m}$.

\begin{trick}[Deleting edges: go backwards]
A DSU cannot split a component. If a problem \emph{removes} roads and asks the same questions, read everything first and
process the removals in reverse order as additions, then print the answers in the original order. If additions and
removals are mixed, see offline dynamic connectivity (Chapter~\ref{ch:dsu}).
\end{trick}
\usetemplate{DSU (union-find)}{Ctrl+Alt+T, U}{cpdsu}

% ======================================================================================================================
\section[Coin combinations]{Coin combinations \hfill \gold}\label{wp:coins}

\begin{statement}{Coin Combinations II}{CSES Problem Set}
Given $n$ distinct coin values and a target $x$, count the ways to make the sum $x$ when the \emph{order of coins does not
matter} ($2 + 3$ and $3 + 2$ are the same way). Each coin value can be used any number of times. Print the count modulo
$10^9+7$. Limits: $n \le 100$, $x \le 10^6$.
\sample{3 9\\2 3 5}{3}
\end{statement}

\wstep{Read the limits.} $n\cdot x = 10^8$ simple steps is fine. The number of ways is astronomically large, hence the
modulus.

\wstep{Find the idea.} Let \cc{ways[s]} be the number of ways to make $s$ with the coins processed so far; start with only
\cc{ways[0] = 1} (the empty sum). Add the coins one at a time: with coin $c$ available, a way to make $s$ either does not use
$c$ at all (the old \cc{ways[s]}), or it is a way to make $s - c$ that may already contain $c$, plus one more $c$. Looping
$s$ \emph{upwards} makes \cc{ways[s - c]} already include uses of $c$, so the coin can be used many times. Putting the coins
in the \emph{outer} loop means every multiset is built in one fixed order (all 2s, then all 3s, then all 5s), so it is
counted exactly once.

\begin{figure}[H]
\centering
\newcommand{\chg}[1]{\cellcolor{cporange!18}\textbf{#1}}
\small
\begin{tabular}{l*{10}{c}}
\toprule
$s$ & 0 & 1 & 2 & 3 & 4 & 5 & 6 & 7 & 8 & 9 \\
\midrule
start & 1 & 0 & 0 & 0 & 0 & 0 & 0 & 0 & 0 & 0 \\
after coin 2 & 1 & 0 & \chg{1} & 0 & \chg{1} & 0 & \chg{1} & 0 & \chg{1} & 0 \\
after coin 3 & 1 & 0 & 1 & \chg{1} & 1 & \chg{1} & \chg{2} & \chg{1} & \chg{2} & \chg{2} \\
after coin 5 & 1 & 0 & 1 & 1 & 1 & \chg{2} & 2 & \chg{2} & \chg{3} & \cellcolor{cpgreen!25}\textbf{3} \\
\bottomrule
\end{tabular}
\caption{\texttt{ways[s]} after each coin; changed entries are highlighted. For example, after coin 5,
\texttt{ways[8]} $= 2 + {}$\texttt{ways[3]} $= 3$: $\{2,2,2,2\}, \{2,3,3\}$ and $\{3,5\}$. The answer is \texttt{ways[9]} $= 3$:
$\{2,2,5\}, \{3,3,3\}, \{2,2,2,3\}$.}
\end{figure}

\wstep{Code.}
\begin{code}
void solve() {
    int n, x;
    cin >> n >> x;
    vi c(n);
    cin >> c;
    vector<ll> ways(x + 1);
    ways[0] = 1;
    for (int coin : c)  // coins in the outer loop: each multiset is counted once
        for (int s = coin; s <= x; s++)
            if ((ways[s] += ways[s - coin]) >= MOD) ways[s] -= MOD;  // cheaper than %
    print(ways[x]);
}
\end{code}

\wstep{Complexity.} $\Oh{n\cdot x}$ time and $\Oh{x}$ memory; 0.21\,s at the maximum size.

\begin{pitfall}[The order of the loops is the problem]
Swapping the loops (sums outside, coins inside) counts \emph{ordered} sequences instead: that is the other CSES problem,
Coin Combinations I. On this sample it gives 8, not 3, because $\{2,2,5\}$ is counted as $2{+}2{+}5$, $2{+}5{+}2$ and
$5{+}2{+}2$, and $\{2,2,2,3\}$ four times.
\end{pitfall}

\begin{tip}
\cc{if ((w += v) >= MOD) w -= MOD;} replaces \cc{w = (w + v) % MOD} whenever both values are already below \cc{MOD}. A
comparison is several times cheaper than a 64-bit division, which matters at $10^8$ steps. Your local build uses the checked
STL (\texttt{-D\_GLIBCXX\_DEBUG}), so this loop runs much slower on your machine than on the judge; GUIDE.pdf,
Section~5.4 shows how to time a solution without it.
\end{tip}
\usetemplate{Knapsack (0/1, unbounded, coin-change ways)}{Ctrl+Alt+T, Shift+Q}{cpknap}

% ======================================================================================================================
\section[Book shop]{Book shop \hfill \gold}\label{wp:books}

\begin{statement}{Book Shop}{CSES Problem Set}
A shop sells $n$ books; book $i$ costs $h_i$ and has $s_i$ pages. With a budget of $x$, buying each book at most once, what
is the largest total number of pages? Limits: $n \le 1000$, $x \le 10^5$, prices and pages up to 1000.
\sample{4 10\\4 8 5 3\\5 12 8 1}{13}
\end{statement}

\wstep{Read the limits.} $2^{1000}$ subsets is impossible, but $n\cdot x = 10^8$ is fine: a DP over (book, budget).

\wstep{Find the idea.} This is the \defn{0/1 knapsack} (Chapter~\ref{ch:dp}). Let \cc{best[c]} be the most pages with total
price at most $c$ using the books seen so far. Adding book $i$ either skips it or buys it:
\[
  \mathrm{best}[c] \leftarrow \max\bigl(\mathrm{best}[c],\ \mathrm{best}[c - h_i] + s_i\bigr).
\]
Iterate $c$ \emph{downwards}, so that \cc{best[c - price[i]]} still describes the books \emph{before} $i$ and book $i$ is bought
at most once. (Upwards is the unbounded version, exactly as in the previous problem.)

\begin{figure}[H]
\centering
\newcommand{\chg}[1]{\cellcolor{cporange!18}\textbf{#1}}
\small
\setlength{\tabcolsep}{4.2pt}
\begin{tabular}{l*{11}{c}}
\toprule
budget $c$ & 0 & 1 & 2 & 3 & 4 & 5 & 6 & 7 & 8 & 9 & 10 \\
\midrule
start & 0 & 0 & 0 & 0 & 0 & 0 & 0 & 0 & 0 & 0 & 0 \\
book 1 (price 4, 5 pages) & 0 & 0 & 0 & 0 & \chg{5} & \chg{5} & \chg{5} & \chg{5} & \chg{5} & \chg{5} & \chg{5} \\
book 2 (price 8, 12 pages) & 0 & 0 & 0 & 0 & 5 & 5 & 5 & 5 & \chg{12} & \chg{12} & \chg{12} \\
book 3 (price 5, 8 pages) & 0 & 0 & 0 & 0 & 5 & \chg{8} & \chg{8} & \chg{8} & 12 & \chg{13} & \chg{13} \\
book 4 (price 3, 1 page) & 0 & 0 & 0 & \chg{1} & 5 & 8 & 8 & 8 & 12 & 13 & \cellcolor{cpgreen!25}\textbf{13} \\
\bottomrule
\end{tabular}
\caption{\texttt{best[c]} after each book. The answer 13 buys books 1 and 3 (price 9).}
\end{figure}

\wstep{Code.}
\begin{code}
void solve() {
    int n, x;
    cin >> n >> x;
    vi price(n), pages(n);
    cin >> price >> pages;
    vector<ll> best(x + 1, 0);  // best[c] = most pages with total price <= c
    rep(i, n)
        for (int c = x; c >= price[i]; c--)  // downwards: every book at most once
            chmax(best[c], best[c - price[i]] + pages[i]);
    print(best[x]);
}
\end{code}

\wstep{Complexity.} $\Oh{n\cdot x}$ time, $\Oh{x}$ memory: one row instead of an $n \times x$ table.

\begin{pitfall}[The greedy by value per price]
Buying the books with the most pages per unit of price is wrong. With budget 10 and books (price 6, 7 pages), (5, 5),
(5, 5): the best ratio takes the first book and then nothing fits (7 pages), but the two others give 10 pages.
\end{pitfall}
\usetemplate{Knapsack (0/1, unbounded, coin-change ways)}{Ctrl+Alt+T, Shift+Q}{cpknap}

% ======================================================================================================================
\section[Longest increasing subsequence]{Longest increasing subsequence \hfill \gold}\label{wp:lis}

\begin{statement}{Increasing Subsequence}{CSES Problem Set}
Given an array of $n$ integers, find the length of its longest strictly increasing subsequence (the elements need not be
adjacent). Limits: $n \le 2\cdot 10^5$, values up to $10^9$.
\sample{8\\7 3 5 3 6 2 9 8}{4}
\end{statement}

\wstep{Read the limits.} The textbook DP, $\mathrm{dp}[i] = 1 + \max\{\mathrm{dp}[j] : j < i,\ a_j < a_i\}$, is
$\Oh{n^2} = 4\cdot10^{10}$. We need $\Oh{n\log n}$.

\wstep{Find the idea.} Keep \cc{tail[k]}: the smallest possible last element of an increasing subsequence of length $k+1$
among the elements seen so far. A small last element is always better, since more future elements can extend it. The
array \cc{tail} is strictly increasing, so for each new $x$ binary search finds the first entry $\ge x$
(\cc{lower_bound}):
\begin{itemize}
  \item if there is none, $x$ extends the longest subsequence: append it;
  \item otherwise a subsequence of that length can now end with the smaller value $x$: overwrite the entry.
\end{itemize}
The length of \cc{tail} at the end is the answer.

\begin{figure}[H]
\centering
\begin{tikzpicture}
  \begin{scope}
    \node[lbl, anchor=east] at (-0.15,0) {$x = 7$}; \arr{T1}{0}{0}{7}\node[cell, hot] at (T1-1.west) {7};
    \node[lbl, anchor=west] at (3,0) {append};
    \node[lbl, anchor=east] at (-0.15,-0.85) {$x = 3$}; \arr{T2}{0}{-0.85}{3}\node[cell, hot] at (T2-1.west) {3};
    \node[lbl, anchor=west] at (3,-0.85) {replaces 7};
    \node[lbl, anchor=east] at (-0.15,-1.7) {$x = 5$}; \arr{T3}{0}{-1.7}{3,5}\node[cell, hot] at (T3-2.west) {5};
    \node[lbl, anchor=west] at (3,-1.7) {append};
    \node[lbl, anchor=east] at (-0.15,-2.55) {$x = 3$}; \arr{T4}{0}{-2.55}{3,5}
    \node[lbl, anchor=west] at (3,-2.55) {3 is there: no change};
  \end{scope}
  \begin{scope}[xshift=7.2cm]
    \node[lbl, anchor=east] at (-0.15,0) {$x = 6$}; \arr{T5}{0}{0}{3,5,6}\node[cell, hot] at (T5-3.west) {6};
    \node[lbl, anchor=west] at (3,0) {append};
    \node[lbl, anchor=east] at (-0.15,-0.85) {$x = 2$}; \arr{T6}{0}{-0.85}{2,5,6}\node[cell, hot] at (T6-1.west) {2};
    \node[lbl, anchor=west] at (3,-0.85) {replaces 3};
    \node[lbl, anchor=east] at (-0.15,-1.7) {$x = 9$}; \arr{T7}{0}{-1.7}{2,5,6,9}\node[cell, hot] at (T7-4.west) {9};
    \node[lbl, anchor=west] at (3,-1.7) {append};
    \node[lbl, anchor=east] at (-0.15,-2.55) {$x = 8$}; \arr{T8}{0}{-2.55}{2,5,6,8}\node[cell, good] at (T8-4.west) {8};
    \node[lbl, anchor=west] at (3,-2.55) {replaces 9; length 4};
  \end{scope}
\end{tikzpicture}
\caption{\texttt{tail} after each element of the sample (read the left column first).}
\end{figure}

\wstep{Code.}
\begin{code}
void solve() {
    int n;
    cin >> n;
    vi a(n);
    cin >> a;
    vi tail;  // tail[k] = smallest end of an increasing subsequence of length k + 1
    for (int x : a) {
        auto it = lower_bound(all(tail), x);
        if (it == tail.end()) tail.push_back(x);
        else *it = x;
    }
    print(sz(tail));
}
\end{code}

\wstep{Complexity.} $\Oh{n\log n}$.

\begin{pitfall}[\texttt{tail} is not a subsequence]
The final \cc{tail} is $2, 5, 6, 8$, but that is not a subsequence of the input (the 2 comes after the 5 and the 6). Only its
\emph{length} is meaningful. To output an actual subsequence, remember for each element its predecessor, as the library
template does. For the longest \emph{non-decreasing} subsequence use \cc{upper_bound} instead.
\end{pitfall}
\usetemplate{Longest increasing subsequence}{Ctrl+Alt+T, Shift+L}{cplis}

% ======================================================================================================================
\section[Subordinates]{Subordinates \hfill \gold}\label{wp:subordinates}

\begin{statement}{Subordinates}{CSES Problem Set}
A company has $n$ employees. Employee 1 is the director and every other employee has one direct boss. For every employee,
count their subordinates (everyone below them in the hierarchy). Limits: $n \le 2\cdot 10^5$.
\sample{5\\1 1 2 3}{4 1 1 0 0}
\end{statement}

\wstep{Read the limits.} Walking up from every employee to the director is $\Oh{n^2}$ on a chain. We need $\Oh{n}$.

\wstep{Find the idea.} The hierarchy is a tree rooted at the director. The subordinates of $u$ are its children and their
subordinates:
\[
  \mathrm{sub}[u] = \sum_{c \text{ child of } u} \bigl(\mathrm{sub}[c] + 1\bigr).
\]
This is a \defn{tree DP} (Chapter~\ref{ch:trees}): every child must be finished before its parent. A recursive DFS does
that, but a chain of $2\cdot10^5$ employees makes it recurse $2\cdot10^5$ deep. The iterative version is just as short: list
the vertices in BFS order from the root (parents before children), then process that list backwards.

\begin{figure}[H]
\centering
\begin{tikzpicture}
  \node[vtx] (1) at (0,0) {1};
  \node[vtx] (2) at (-1.2,-1.2) {2}; \node[vtx] (3) at (1.2,-1.2) {3};
  \node[vtx] (4) at (-1.2,-2.4) {4}; \node[vtx] (5) at (1.2,-2.4) {5};
  \draw[ed] (1) -- (2) (1) -- (3) (2) -- (4) (3) -- (5);
  \node[lbl, right=3pt of 1, text=cporange] {sub $= 4$};
  \node[lbl, left=3pt of 2, text=cporange] {sub $= 1$}; \node[lbl, right=3pt of 3, text=cporange] {sub $= 1$};
  \node[lbl, left=3pt of 4, text=cporange] {sub $= 0$}; \node[lbl, right=3pt of 5, text=cporange] {sub $= 0$};
  \arr{O}{4}{-1.2}{1,2,3,4,5}
  \node[lbl, above=2pt of O-3] {BFS order (0-indexed in the code)};
  \draw[hded] ($(O-5.south)+(0,-0.35)$) -- node[below, lbl] {processed right to left: children first} ($(O-1.south)+(0,-0.35)$);
\end{tikzpicture}
\caption{The sample hierarchy with its answers, and the order in which the DP visits it.}
\end{figure}

\wstep{Code.}
\begin{code}
void solve() {
    int n;
    cin >> n;
    vvi ch(n);
    rep(i, 1, n) {
        int b;
        cin >> b;
        ch[b - 1].push_back(i);  // employee i's boss is b
    }
    vi order{0};  // BFS order from the director: parents before children
    rep(i, n) for (int c : ch[order[i]]) order.push_back(c);
    vi sub(n, 0);
    per(i, n) {  // children before parents
        int u = order[i];
        for (int c : ch[u]) sub[u] += sub[c] + 1;
    }
    print(sub);
}
\end{code}

\wstep{Complexity.} $\Oh{n}$.

\begin{tip}
\cc{order} grows while the loop reads it. That is safe with an index, as here, but \emph{not} with a range-based \cc{for}
or an iterator: \cc{push_back} may move the vector's storage and leave the iterator dangling.
\end{tip}
\usetemplate{Tree DP (subtree DP, iterative)}{Ctrl+Alt+T, ;}{cptreedp}

% ======================================================================================================================
\section[Tree diameter]{Tree diameter \hfill \gold}\label{wp:diameter}

\begin{statement}{Tree Diameter}{CSES Problem Set}
Given a tree with $n$ nodes, find its diameter: the largest number of edges on a path between two nodes.
Limits: $n \le 2\cdot 10^5$.
\sample{5\\1 2\\1 3\\3 4\\3 5}{3}
\end{statement}

\wstep{Read the limits.} A BFS from every node is $\Oh{n^2}$. We need $\Oh{n}$.

\wstep{Find the idea.} Run a BFS from any node; the node \cc{far} that is farthest from it is always one end of some
diameter. Run a second BFS from \cc{far}: the largest distance found is the diameter. (Sketch: if a longest path did not end
at \cc{far}, you could splice the path to \cc{far} onto it and make it longer; Chapter~\ref{ch:trees} has the details and
the DP alternative.)

\begin{figure}[H]
\centering
\begin{tikzpicture}
  \begin{scope}
    \node[vtx, hot] (1) at (0,0) {1}; \node[vtx] (2) at (-1,-1.1) {2}; \node[vtx] (3) at (1,-1.1) {3};
    \node[vtx, good] (4) at (0.4,-2.3) {4}; \node[vtx] (5) at (1.6,-2.3) {5};
    \draw[ed] (1) -- (2) (1) -- (3) (3) -- (4) (3) -- (5);
    \node[lbl, above=2pt of 1] {0}; \node[lbl, left=2pt of 2] {1}; \node[lbl, right=2pt of 3] {1};
    \node[lbl, below=2pt of 4] {2}; \node[lbl, below=2pt of 5] {2};
    \node[note, text width=4.4cm] at (0.3,-3.85) {BFS from 1: the farthest node is 4 (the first with distance 2)};
  \end{scope}
  \begin{scope}[xshift=6.5cm]
    \node[vtx, good] (1) at (0,0) {1}; \node[vtx, good] (2) at (-1,-1.1) {2}; \node[vtx, good] (3) at (1,-1.1) {3};
    \node[vtx, hot] (4) at (0.4,-2.3) {4}; \node[vtx] (5) at (1.6,-2.3) {5};
    \draw[ed] (3) -- (5);
    \draw[hed] (4) -- (3) -- (1) -- (2);
    \node[lbl, above=2pt of 1] {2}; \node[lbl, left=2pt of 2] {3}; \node[lbl, right=2pt of 3] {1};
    \node[lbl, below=2pt of 4] {0}; \node[lbl, below=2pt of 5] {2};
    \node[note, text width=4.4cm] at (0.3,-3.85) {BFS from 4: the largest distance is 3, along the diameter 4--3--1--2};
  \end{scope}
\end{tikzpicture}
\caption{The two BFS passes on the sample; numbers next to the nodes are distances.}
\end{figure}

\wstep{Code.}
\begin{code}
void solve() {
    int n;
    cin >> n;
    vvi adj(n);
    rep(i, n - 1) {
        int a, b;
        cin >> a >> b;
        a--, b--;
        adj[a].push_back(b), adj[b].push_back(a);
    }
    auto bfs = [&](int s) {  // distances from s
        vi d(n, -1);
        queue<int> q;
        q.push(s), d[s] = 0;
        while (!q.empty()) {
            int u = q.front();
            q.pop();
            for (int v : adj[u])
                if (d[v] < 0) d[v] = d[u] + 1, q.push(v);
        }
        return d;
    };
    vi d0 = bfs(0);
    int far = max_element(all(d0)) - d0.begin();  // one end of a diameter
    vi df = bfs(far);
    print(*max_element(all(df)));
}
\end{code}

\wstep{Complexity.} $\Oh{n}$: two BFS passes. With $n = 1$ there are no edges and the answer is 0, which the code handles.

\begin{pitfall}
The two-pass trick needs a \emph{tree} with non-negative edge weights (with weights, use a DFS that adds them up instead of a
BFS). On general graphs the longest simple path is NP-hard, and the trick gives wrong answers.
\end{pitfall}
\usetemplate{Tree diameter}{Ctrl+Alt+T, Shift+D}{cpdiameter}

% ======================================================================================================================
\section[Range sum queries with updates]{Range sum queries with updates \hfill \gold}\label{wp:bit}

\begin{statement}{Dynamic Range Sum Queries}{CSES Problem Set}
Given an array of $n$ integers, process $q$ queries: \texttt{1 k u} sets $a_k = u$; \texttt{2 a b} prints
$a_a + \dots + a_b$ (positions are 1-indexed). Limits: $n, q \le 2\cdot 10^5$, values up to $10^9$.
\sample{8 4\\3 2 4 5 1 1 5 3\\2 1 4\\2 5 6\\1 3 1\\2 1 4}{14\\2\\11}
\end{statement}

\wstep{Read the limits.} Prefix sums answer a sum in $\Oh{1}$ but an update changes $\Oh{n}$ of them; a plain array updates
in $\Oh{1}$ but sums in $\Oh{n}$. Either way the worst case is $\Oh{nq} = 4\cdot10^{10}$. Both operations must be about
$\Oh{\log n}$.

\wstep{Find the idea.} A \defn{Fenwick tree} (binary indexed tree, Chapter~\ref{ch:segtree}) does both in $\Oh{\log n}$.
Entry $t_i$ stores the sum of the $\mathrm{lowbit}(i)$ positions that end at $i$, where $\mathrm{lowbit}(i)$ is the lowest
set bit of $i$. A prefix sum adds $\Oh{\log n}$ entries, removing the lowest bit each time; an update changes
$\Oh{\log n}$ entries, adding the lowest bit each time. The tree supports ``add $v$ at $i$'', so ``set $a_k = u$'' becomes
``add $u - a_k$'', and we keep the current array to know $a_k$.

\begin{figure}[H]
\centering
\begin{tikzpicture}
  \foreach \v [count=\i] in {3,2,4,5,1,1,5,3} {
    \node[cell, minimum width=8mm] at ({(\i-1)*0.8}, 0) {\v};
    \node[lbl] at ({(\i-0.5)*0.8}, -0.55) {\i};
  }
  \foreach \i/\lb/\lv/\v/\st in {1/1/0/3/seg, 2/2/1/5/seg, 3/1/0/4/seg, 4/4/2/14/hot, 5/1/0/1/hot, 6/2/1/2/seg, 7/1/0/5/cool,
                                 8/8/3/24/cool}
    \node[seg, \st, minimum width={\lb*0.8cm - 0.06cm}, anchor=west] at ({(\i-\lb)*0.8 + 0.03}, {0.85 + \lv*0.72})
      {$t_{\i} = \v$};
  \node[note, text width=5.6cm, anchor=west] at (6.8, 1.5)
    {Each bar is one entry $t_i$ and the positions it sums. Orange: prefix$(5) = t_5 + t_4 = 1 + 14 = 15$ (index
     $5 \to 4 \to 0$). Teal: changing position 7 updates $t_7$ and $t_8$ (index $7 \to 8$).};
\end{tikzpicture}
\caption{A Fenwick tree over the sample array.}
\end{figure}

\wstep{Code.} Paste the Fenwick template above \cc{solve()} with \key{Ctrl+Alt+T, I}. It is 0-indexed with half-open ranges
$[l, r)$, so positions $a..b$ of the statement are \cc{sum(a - 1, b)}:
\begin{code}
void solve() {  // BIT (cpbit) is pasted above solve()
    int n, q;
    cin >> n >> q;
    vi a(n);
    cin >> a;
    BIT bit(n);
    rep(i, n) bit.add(i, a[i]);
    while (q--) {
        int type, x, y;
        cin >> type >> x >> y;
        if (type == 1) x--, bit.add(x, y - a[x]), a[x] = y;  // set a[x] = y
        else print(bit.sum(x - 1, y));  // positions x..y, 1-indexed
    }
}
\end{code}

\wstep{Complexity.} $\Oh{(n + q)\log n}$.

\begin{recognize}
Point updates and range sums (or any operation that can be undone, like XOR): Fenwick tree. Range minimum or maximum with
updates cannot be undone by subtraction: use a segment tree (Chapter~\ref{ch:segtree}). No updates at all: plain prefix sums.
\end{recognize}
\usetemplate{Fenwick tree / BIT}{Ctrl+Alt+T, I}{cpbit}

% ======================================================================================================================
\section[String matching]{String matching \hfill \gold}\label{wp:kmp}

\begin{statement}{String Matching}{CSES Problem Set}
Given a string and a pattern, count the positions where the pattern occurs in the string (occurrences may overlap).
Limits: both lengths up to $10^6$.
\sample{saippuakauppias\\pp}{2}
\end{statement}

\wstep{Read the limits.} Comparing the pattern at every position costs $\Oh{|t|\cdot|p|}$, up to $10^{12}$ on inputs like
\texttt{aaa\ldots a} and \texttt{aa\ldots ab}. We need $\Oh{|t| + |p|}$.

\wstep{Find the idea.} The \defn{prefix function} $\pi[i]$ of a string is the length of the longest proper prefix of
$s[0..i]$ that is also a suffix of it (Chapter~\ref{ch:strings}). Compute it on $s = p + \texttt{\#} + t$. The separator
appears in neither string, so no border can be longer than $|p|$, and $\pi[i] = |p|$ exactly when the pattern ends at
position $i$. Computing $\pi$ takes $\Oh{|s|}$ in total: $j$ grows by at most one per position and every step of the inner
\cc{while} loop shrinks it.

\begin{figure}[H]
\centering
\begin{tikzpicture}
  \arr{S}{0}{0}{a,b,a,\#,a,b,a,b,a}\node[lbl, left=4pt of S-1] {$s$};
  \arr{P}{0}{-0.9}{0,0,1,0,1,2,3,2,3}\node[lbl, left=4pt of P-1] {$\pi$};\idx{P}{9}
  \node[cell, good] at (P-7.west) {3}; \node[cell, good] at (P-9.west) {3};
  \draw[brace] (S-1.north west) -- node[above=7pt, lbl] {pattern} (S-3.north east);
  \draw[brace] (S-5.north west) -- node[above=7pt, xshift=-4mm, lbl] {match 1} (S-7.north east);
  \draw[brace, decoration={raise=17pt}] (S-7.north west) -- node[above=22pt, lbl] {match 2} (S-9.north east);
  \node[note, text width=4.6cm, anchor=west] at (6.8, -0.45)
    {Pattern \texttt{aba}, text \texttt{ababa}. $\pi = 3 = |p|$ (green) marks the end of an occurrence; the two
     occurrences overlap.};
\end{tikzpicture}
\caption{The prefix function of pattern\,\texttt{\#}\,text.}
\end{figure}

\wstep{Code.}
\begin{code}
void solve() {
    string text, pat;
    cin >> text >> pat;
    string s = pat + '#' + text;
    int m = sz(pat), cnt = 0;
    vi pi(sz(s));
    rep(i, 1, sz(s)) {  // prefix function of pat#text
        int j = pi[i - 1];
        while (j > 0 && s[i] != s[j]) j = pi[j - 1];
        if (s[i] == s[j]) j++;
        pi[i] = j;
        cnt += j == m;  // a border of length |pat| ends here: one occurrence
    }
    print(cnt);
}
\end{code}

\wstep{Complexity.} $\Oh{|t| + |p|}$ time and memory; 0.05\,s at the maximum size.

\begin{pitfall}
The separator must not occur in either string. CSES strings are lowercase letters, so \texttt{\#} is safe; for arbitrary
bytes, run the prefix function on the pattern alone and match the text against it (the library's \cc{find_occurrences}).
\end{pitfall}
\usetemplate{KMP (prefix function + pattern search)}{Ctrl+Alt+T, Shift+K}{cpkmp}

% ======================================================================================================================
\section[Teleporters]{Teleporters: binary lifting \hfill \gold}\label{wp:planets}

\begin{statement}{Planets Queries I}{CSES Problem Set}
There are $n$ planets and each has one teleporter, which sends you to planet $t_i$. Answer $q$ queries: starting from planet
$x$ and using a teleporter $k$ times, where do you end up? Limits: $n, q \le 2\cdot10^5$, $k \le 10^9$.
\sample{4 3\\2 1 1 4\\1 2\\3 4\\4 1}{1\\2\\4}
\end{statement}

\wstep{Read the limits.} Simulating $k \le 10^9$ steps per query is far too slow. Each query must take about $\Oh{\log k}$.

\wstep{Find the idea.} Every planet has exactly one outgoing edge: a \defn{functional graph} (Chapter~\ref{ch:graphs}).
Precompute \cc{up[j][v]}, the planet reached from $v$ after $2^j$ steps. Two jumps of $2^{j-1}$ make one jump of $2^j$:
\[
  \mathrm{up}[j][v] = \mathrm{up}[j-1]\bigl[\mathrm{up}[j-1][v]\bigr].
\]
A query writes $k$ in binary and makes one jump for every set bit, at most 30 jumps. This is \defn{binary lifting}, the same
table that finds ancestors and LCAs in trees (Chapter~\ref{ch:trees}).

\begin{figure}[H]
\centering
\begin{minipage}{0.34\linewidth}
\centering
\begin{tikzpicture}
  \node[vtx] (1) at (0,0) {1}; \node[vtx] (2) at (1.9,0) {2};
  \node[vtx] (3) at (0,-1.5) {3}; \node[vtx] (4) at (1.9,-1.5) {4};
  \draw[ded] (1) to[bend left=25] (2); \draw[ded] (2) to[bend left=25] (1);
  \draw[ded] (3) -- (1);
  \draw[ded] (4) to[out=25, in=-25, looseness=7] (4);
\end{tikzpicture}
\end{minipage}\hfill
\begin{minipage}{0.62\linewidth}
\small
\begin{tabular}{@{}lcccc@{}}
\toprule
 & $v = 1$ & $v = 2$ & $v = 3$ & $v = 4$ \\
\midrule
\texttt{up[0][v]} (1 step) & 2 & 1 & 1 & 4 \\
\texttt{up[1][v]} (2 steps) & 1 & 2 & 2 & 4 \\
\texttt{up[2][v]} (4 steps) & 1 & 2 & \cellcolor{cporange!25}\textbf{2} & 4 \\
\bottomrule
\end{tabular}

\smallskip
{\footnotesize Query $(x = 3, k = 4)$: $4 = 100_2$, so one jump: \texttt{up[2][3]} $= 2$.}
\end{minipage}
\caption{The functional graph of the sample and the first rows of its jump table.}
\end{figure}

\wstep{Code.}
\begin{code}
void solve() {
    int n, q;
    cin >> n >> q;
    const int LOG = 30;  // k < 2^30
    vector<vi> up(LOG, vi(n));  // up[j][v] = planet reached from v after 2^j teleports
    rep(i, n) cin >> up[0][i], up[0][i]--;
    rep(j, 1, LOG) rep(i, n) up[j][i] = up[j - 1][up[j - 1][i]];
    while (q--) {
        int x, k;
        cin >> x >> k;
        x--;
        rep(j, LOG) if (k >> j & 1) x = up[j][x];
        print(x + 1);
    }
}
\end{code}

\wstep{Complexity.} $\Oh{n\log K}$ to build the table and $\Oh{\log K}$ per query, with $K = 10^9$; 0.21\,s at the maximum
size.

\begin{tip}[Memory]
The table holds $30n$ numbers: 48\,MB with the template's 64-bit \cc{int}. That is fine for CSES (512\,MB) but on a 256\,MB
problem with bigger $n$, store it as \cc{vector<vector<int32_t>>} to halve it.
\end{tip}
\usetemplate{K-th successor (binary lifting / jumping)}{Ctrl+Alt+T, .}{cpjump}

% ======================================================================================================================
\section[Where to go next]{Where to go next \hfill \bronze}

\begin{summary}[What these problems teach]
\begin{itemize}
  \item \textbf{Ask a smaller question}: ``best subarray \emph{ending here}'' (Kadane), ``smallest tail for each length''
        (LIS), ``can it be done in time $T$?'' (binary search on the answer).
  \item \textbf{Change the representation}: subarray sums become pairs of prefix sums, a grid becomes a graph, a hierarchy
        becomes a tree.
  \item \textbf{Prove the greedy} with an exchange argument, or break it with a small counterexample.
  \item \textbf{Loop order is meaning}: coins outside or inside, budget upwards or downwards.
  \item \textbf{Check every bound}: 64-bit sums, overflow inside a check, recursion depth, memory of a table.
\end{itemize}
\end{summary}

The CSES Problem Set has about 300 problems in the same style, grouped by topic, and it is the best next step after this
chapter; the USACO Guide lists problems for every section of this book by level. Solve a few from each topic, and for every
problem you could not solve, read the editorial and implement it yourself the same day.
